This thesis has developed a procedural music generator that turns explicit music-theory principles into a working software system. Over three iterations, the project moved from a small proof-of-concept script to a modular generator that can create melodies and harmonized outputs, render them as notation and audio, and be steered through concrete musical choices. The final iteration is therefore not just a better prototype, but a finished system with a clear musical model and a clear technical structure.

The main findings of the thesis are:
\begin{itemize}
  \item procedural generation can create structured musical material without relying on large training data sets
  \item explicit representations of motifs, form sections, harmony spans, and voice profiles give meaningful control over the generated music
  \item a rule-based generator can be guided by key, rhythm, form, texture, and harmony while still remaining transparent to the user
  \item iterative development was the right way to solve the problem, because each iteration exposed the weaknesses the next one had to address
\end{itemize}

The final system does not merely produce notes that stay inside a key. It produces musical examples with phrase shape, motivic reuse, harmonic direction, range awareness, and controllable output formats. It can be used as a compositional aid, as a rapid prototyping tool for melodic and harmonic ideas, and as a technical foundation for continued work in procedural or mixed-initiative music generation. The thesis therefore contributes both a generator and a practical architectural approach to building theory-aware music software.

The use of \ac{AI} in development also led to a clear conclusion. Codex was useful when it was used as a tool for implementation, refactoring, debugging, and iteration inside an already understood problem space. It accelerated development, but it did not replace musical judgment, software design judgment, or responsibility for the result. That is an important finding, because it places \ac{AI} in a realistic and useful role inside software engineering work.

Within the scope of the thesis, the problem has been solved. The project set out to design a system capable of generating coherent musical material from explicit musical principles while also examining iterative development and \ac{AI}-assisted software development. The final system achieves the goal and the thesis evaluates iterative development and \ac{AI} as a development tool. It generates structured musical examples and supports a wide range of controlled inputs.
